\Needspace{10\baselineskip}
\subsection{The whole effective definable concordance model}

\begin{saresult}{A.161}{Proved}{Whole definable concordance spaces from finite algebraic codes}
Let \(M\) be any fixed o-minimal expansion of a real closed field,
let \(B\) be a closed definably compact definable manifold, and choose
a definable triangulation \(\tau:|K|_M\to B\) with \(K\) a finite
abstract PL manifold. Let \(G=G^{\rm PL}(K)\) be the explicitly coded
simplicial group of A.131. Then there is a specified weak equivalence
\[
 G\longrightarrow P_{\rm def,M}(B).
 \tag{161.1}
\]
It is the composite of real-algebraic coefficient extension, PL
inclusion, and conjugation by \(\tau\times\operatorname{id}_I\).

In particular, for every \(r\geq0\) the complete recursive presentation
\(\Pi_r^{\rm PL}(K)\) of A.131 presents
\(\pi_rP_{\rm def,M}(B)\). For \(r=0\) this computes the full component
group. For positive \(r\) it is based at the identity, with the other
components obtained by translation. For the real realization
\(|K|_{\mathbb R}\), the same effective model also computes TOP
concordances whenever every component has dimension outside three
and four.

The algorithm takes only the finite complex \(K\), not a decision
procedure or an oracle for the elementary theory of \(M\).
The output is the complete countable simplicial model and complete
recursively enumerable group presentations. It does not assert a
terminating word problem or a finite invariant-factor decomposition.
\end{saresult}

\begin{proof}
We recall the finite nature of the already proved construction,
so that the resulting whole-space answer is explicit. Realize \(K\)
by standard rational vertices and put \(F=\mathbb R_{\rm alg}\).
A degree-\(q\) code is the least binary string describing a finite
union of simplices with real-algebraic vertices in coordinates
\((z,z',t)\), where \(z,z'\in |K|\times I\) and
\(t\in\Delta^q\). The parameter variable is literally shared.
The two graph projections must be bijective; the bottom must be fixed
pointwise and the free top preserved setwise. Every one of these
conditions, and equality of two such graphs, is a finite linear
first-order test over the effectively ordered field \(F\).
A.131 proves that precisely these codes describe jointly PL
parameter-preserving homeomorphisms, including joint continuity of
the map and inverse.

The same proof gives terminating algorithms for products, inverses,
faces, degeneracies, and canonicalization. Thus the resulting
simplicial group is literally isomorphic to
\(P_{\rm PL,F}(|K|_F)\). A.115 gives a weak equivalence
\[
 G=P_{\rm PL,F}(|K|_F)
 \longrightarrow P_{\rm PL,M}(|K|_M).
\]
A.160 gives a weak equivalence from the latter group to
\(P_{\rm def,M}(|K|_M)\). Conjugation by the chosen definable
triangulation is a literal isomorphism of the final definable
simplicial group with \(P_{\rm def,M}(B)\): on a family \(h_p\) it is
\[
 h_p\longmapsto
 (\tau\times\operatorname{id}_I)\,h_p\,
 (\tau\times\operatorname{id}_I)^{-1}.
\]
The formula preserves the parameter, the bottom identity, and the
top marking. Its inverse is conjugation by the inverse triangulation.
Composing proves (161.1). The triangulation is part of the specified
comparison data; no literal map from non-Archimedean points to real
points is asserted.

For completeness, the presentation is as follows. In the effective
simplicial group \(G\), put
\[
 \begin{aligned}
 Z_0&=G_0,\\
 Z_r&=\bigcap_{i=0}^{r}\ker(d_i:G_r\to G_{r-1})\quad(r\geq1),\\
 N_{r+1}&=\bigcap_{i=1}^{r+1}\ker(d_i:G_{r+1}\to G_r).
 \end{aligned}
\]
All these sets have decidable membership on finite codes. The Moore
quotient theorem, in the exact form and convention proved in A.131,
identifies \(\pi_r|G|=Z_r/d_0N_{r+1}\). Consequently the full group is
\[
 \left\langle x_a\ (a\in Z_r)\ \left|\
 \begin{array}{l}
 x_1=1,\quad x_ax_b=x_{ab}\quad(a,b\in Z_r),\\
 x_{d_0u}=1\quad(u\in N_{r+1})
 \end{array}\right.\right\rangle .
 \tag{161.2}
\]
Enumerating canonical codes and applying the finite tests enumerates
every generator and every relation in (161.2). The multiplication
relations present the entire group \(Z_r\); the last relations
quotient by precisely the Moore boundary subgroup, not an unspecified
subset of nullhomotopies. This proves the component statement first,
with \(N_1=\ker d_1\), and then every higher group. For \(r\geq1\)
the quotient is abelian, so A.131's integer relation matrix is an
equivalent output.

For a real base with all component dimensions at most two or at least
five, A.134 and the high-dimensional comparison used in A.113 give
the PL-to-TOP weak equivalence. Applying it componentwise gives the
TOP assertion. An empty base gives the constant one-element group.
There are only finitely many components; a bottom-fixed concordance
preserves each corresponding product component, so this
componentwise treatment is exact.

A.131's graph-code validity, operations, and Moore presentation were
proved and independently checked before the present comparison.
We have not reinterpreted its finite check as a verification of
A.160. The new logical input here is that A.160 now permits the
entire effective model to be used in every fixed o-minimal expansion,
without enumerating formulas of that expansion. General real TOP
bases of dimensions three and four remain outside this conclusion.
\end{proof}

